\chapter{Anti-CCP Antibody (Cyclic Citrullinated Peptide)}

\section*{What Is This Test?}
The anti-cyclic citrullinated peptide (anti-CCP) antibody test detects autoantibodies directed against citrullinated proteins. Citrullination is a post-translational modification in which the amino acid arginine is converted to citrulline by the enzyme peptidylarginine deiminase (PAD). In rheumatoid arthritis, aberrant citrullination of synovial proteins (including filaggrin, vimentin, fibrinogen, and type II collagen) generates neoepitopes that trigger an immune response, leading to the production of anti-CCP antibodies. These antibodies are highly specific for rheumatoid arthritis (specificity 95--98\%) and can be detected years before the onset of clinical symptoms \cite{vanvenrooij2011}. Anti-CCP positivity at diagnosis is a strong predictor of erosive, aggressive disease. Testing is performed by enzyme-linked immunosorbent assay (ELISA) using synthetic cyclic citrullinated peptides (second-generation CCP2 assay) or by automated chemiluminescent immunoassays.

\section*{Alternative Names}
Anti-CCP; Anti-Cyclic Citrullinated Peptide Antibody; CCP Antibody; ACPA (Anti-Citrullinated Protein Antibody); Anti-CCP2; Citrullinated Peptide Antibody

\section*{Principle}
Anti-CCP antibodies are measured using second-generation enzyme-linked immunosorbent assay (CCP2 ELISA) or automated chemiluminescent immunoassays (CLIA). Synthetic cyclic citrullinated peptides, designed to mimic the conformational epitopes of citrullinated proteins found in the rheumatoid joint, are coated onto microtiter plates or paramagnetic microparticles. Diluted patient serum is incubated with the antigen, allowing anti-CCP antibodies to bind. After washing, an enzyme-conjugated (horseradish peroxidase or alkaline phosphatase) anti-human IgG detection antibody is added. Following a second wash, a chemiluminescent substrate or chromogenic substrate (tetramethylbenzidine, TMB) is added, and the resulting signal is measured photometrically or luminometrically. The signal intensity is proportional to the anti-CCP antibody concentration. Results are reported in arbitrary units (U/mL) with calibration against a reference standard. The CCP2 assay (second generation) has superior sensitivity compared to the first-generation CCP1 assay, with the CCP3 and CCP3.1 assays offering further refinements in epitope targeting.

\section*{History}
The discovery of anti-perinuclear factor (APF) in 1964 and anti-keratin antibody (AKA) in 1979 provided the first evidence for autoantibodies targeting citrullinated proteins in RA, though these tests were cumbersome (requiring human buccal mucosa or rat esophagus as substrate). In 1998, Walther van Venrooij and colleagues at the University of Nijmegen demonstrated that citrullinated filaggrin was the target antigen of APF and AKA \cite{schellekens1998}. They developed the first-generation CCP ELISA (CCP1) in 2000 using linear citrullinated peptides. The CCP2 assay, introduced in 2002, used cyclic peptides that better mimicked the native conformational epitopes, dramatically improving sensitivity (from 55\% to 75\%) while maintaining >95\% specificity. The recognition that anti-CCP antibodies can precede clinical RA by 5--10 years was a landmark finding reported by Aho et al. and Rantap\"a\"a-Dahlqvist et al. in the early 2000s. The 2010 ACR/EULAR classification criteria for RA formally incorporated anti-CCP as a serologic criterion alongside rheumatoid factor, reflecting its importance in RA diagnosis \cite{aletaha2010}.

\section*{Why This Test Is Ordered}
\begin{itemize}
  \item Diagnosis of rheumatoid arthritis in patients with early inflammatory arthritis --- anti-CCP has 95--98\% specificity for RA, significantly higher than rheumatoid factor (75--80\%)
  \item Classification of RA according to the 2010 ACR/EULAR criteria, where anti-CCP positivity contributes to the serology score
  \item Prognostic assessment at the time of RA diagnosis --- anti-CCP positivity predicts more aggressive, erosive disease with radiographic progression
  \item Evaluation of patients with undifferentiated inflammatory arthritis to predict progression to RA
  \item Differentiation of RA from other causes of inflammatory arthritis (psoriatic arthritis, reactive arthritis, gout, pseudogout)
  \item Risk stratification before initiating disease-modifying antirheumatic drug (DMARD) therapy, as anti-CCP-positive patients typically require more intensive treatment to prevent joint damage
  \item Testing of asymptomatic individuals with a family history of RA in research settings, as anti-CCP can be present years before clinical onset
\end{itemize}

\section*{Primary vs Secondary Test}
This is a primary diagnostic test for rheumatoid arthritis. Anti-CCP testing is recommended alongside rheumatoid factor in the initial evaluation of suspected RA. Because of its superior specificity, anti-CCP is increasingly ordered as a first-line test in preference to rheumatoid factor.

\section*{How This Test Is Performed}
For most patients, a health care professional will take a blood sample from a vein in your arm using a small needle. After the needle is inserted, a small amount of blood will be collected into a test tube or vial. You may feel a little sting when the needle goes in or out. This usually takes less than five minutes.

\section*{What Preparation Is Needed}
No special preparation is required. Fasting is not necessary. Patients should inform their healthcare provider about all medications. Biologic and targeted synthetic DMARDs may lower anti-CCP levels, though antibodies rarely become completely negative with treatment.

\section*{Contraindications and Precautions}
No absolute contraindications. Specimen should be collected in a serum separator tube (SST, gold-top) or red-top tube. Grossly hemolyzed or lipemic specimens may interfere with ELISA measurements. Anti-CCP positivity is highly specific for RA but can rarely be found in other conditions: psoriatic arthritis (5--10\%), active tuberculosis, chronic hepatitis C infection, and Sjogren syndrome. Unlike rheumatoid factor, anti-CCP positivity does not increase with age and is uncommon in healthy individuals (<1--2\%). Anti-CCP antibodies are relatively stable, and serial testing for monitoring disease activity is not recommended --- clinical assessment and acute phase reactants (CRP, ESR) are preferred for monitoring.

\section*{Risks}
Minimal risk. Slight pain or bruising at the venipuncture site. Rare: hematoma, infection, vasovagal syncope.

\section*{What Happens After the Test}
Results are typically available within 1 to 3 days. A positive anti-CCP result in a patient with symmetric polyarthritis strongly supports the diagnosis of RA and identifies patients at risk for erosive disease. A negative anti-CCP does not exclude RA --- approximately 25--35\% of patients with established RA are anti-CCP-negative (seronegative RA). If RA is strongly suspected despite negative anti-CCP and negative RF, the diagnosis is made on clinical and radiographic grounds. Your healthcare provider will interpret results in the context of your symptoms, joint examination, and inflammatory markers.

\section*{Understanding Results}

\subsection*{Below Normal (Negative)}
\begin{itemize}
  \item Seronegative rheumatoid arthritis: 25--35\% of patients with established RA test negative for anti-CCP. Seronegative RA is generally associated with a less aggressive disease course with less radiographic progression
  \item Early rheumatoid arthritis: A small proportion of patients with early RA are anti-CCP-negative at presentation but may seroconvert over time, though persistent seronegativity is more common
  \item Non-RA inflammatory arthritis: Anti-CCP is negative in psoriatic arthritis (90--95\%), reactive arthritis, ankylosing spondylitis, gout, pseudogout, and osteoarthritis
  \item Viral polyarthritis (parvovirus B19, hepatitis B, hepatitis C, rubella, chikungunya) is typically anti-CCP-negative
  \item Systemic lupus erythematosus with inflammatory arthritis is typically anti-CCP-negative, distinguishing it from RA overlap
\end{itemize}

\subsection*{Above Normal (Positive)}
\begin{itemize}
  \item Rheumatoid arthritis: Anti-CCP is positive in 65--75\% of patients with RA. High-titer anti-CCP (>3 times upper limit of normal) strongly predicts erosive, progressive disease with worse functional outcomes
  \item Pre-clinical RA: Anti-CCP antibodies can be detected 5--10 years before the onset of clinical synovitis. Approximately 40--50\% of anti-CCP-positive individuals with arthralgias but no clinical synovitis will develop RA within 3--5 years
  \item Psoriatic arthritis: Anti-CCP is positive in 5--10\% of patients with psoriatic arthritis, often those with polyarticular disease resembling RA
  \item Sjogren syndrome with articular involvement: Anti-CCP is positive in a small subset (<5\%) of patients
  \item Active pulmonary tuberculosis: Anti-CCP positivity has been reported in 20--30\% of patients with active TB due to citrullination of lung proteins during granulomatous inflammation
  \item Anti-CCP titer does not typically correlate with disease activity and should not be used to monitor treatment response
\end{itemize}

\section*{Normal Results}
\begin{itemize}
  \item \textbf{Anti-CCP (CCP2 assay):} <20 U/mL (negative)
  \item \textbf{Weak positive:} 20--39 U/mL
  \item \textbf{Moderate positive:} 40--60 U/mL
  \item \textbf{Strong positive:} >60 U/mL (>3 times upper limit of normal); this level contributes 3 points to the 2010 ACR/EULAR RA classification criteria serology score
  \item \textbf{Healthy population:} <1--2\% false-positive rate, compared to 5--25\% for rheumatoid factor
  \item \textbf{Elderly (>65 years):} False-positive rate remains <2\%, unlike RF which increases with age
  \item Anti-CCP results should always be interpreted alongside rheumatoid factor, clinical symptoms, and inflammatory markers (CRP, ESR)
\end{itemize}

\section*{Follow-up Tests}
\begin{itemize}
  \item \textbf{Rheumatoid factor (RF):} Combined testing of anti-CCP and RF increases sensitivity for RA diagnosis. Double positivity (anti-CCP-positive and RF-positive) confers the highest risk for aggressive erosive disease
  \item \textbf{C-reactive protein (CRP) and erythrocyte sedimentation rate (ESR):} Acute phase reactants used to assess systemic inflammation and monitor treatment response in RA
  \item \textbf{Antinuclear antibody (ANA):} If features of systemic lupus erythematosus or other connective tissue disease are present alongside arthritis
  \item \textbf{Radiographs of hands, wrists, and feet:} Baseline radiographs to assess for existing erosive changes, joint space narrowing, and to provide a reference for monitoring disease progression
  \item \textbf{Musculoskeletal ultrasound or MRI:} To detect early synovitis, tenosynovitis, and bone erosions not yet visible on conventional radiographs
  \item \textbf{Complete blood count:} To evaluate for anemia of chronic disease, which is common in active RA
  \item \textbf{Comprehensive metabolic panel:} Baseline assessment of renal and hepatic function prior to initiating DMARD therapy
  \item \textbf{Hepatitis B and C serology, latent tuberculosis screening (IGRA or PPD):} Before initiating biologic DMARD therapy, per treatment guidelines
  \item \textbf{HLA typing (HLA-DRB1 shared epitope):} Research tool not routinely used in clinical practice; the presence of the shared epitope alleles is associated with anti-CCP positivity and more severe disease
\end{itemize}

\section*{References}
\bibliographystyle{unsrtnat}
\bibliography{references}

\clearpage
\section*{Regulatory Status and Authorization Date}
Regulatory status applies to the specific assay, instrument, manufacturer, and intended use rather than to the test name alone. No single approval date can be assigned to this test category. For a commercial assay, verify the current product in the relevant regulator's database: the U.S. Food and Drug Administration (FDA) clearance, approval, or authorization; India's Central Drugs Standard Control Organisation (CDSCO) licence or permission; the European Union In Vitro Diagnostic Regulation (EU IVDR) CE certificate issued through the applicable conformity-assessment route; or the United Kingdom Medicines and Healthcare products Regulatory Agency (MHRA) registration and UKCA/CE status. Laboratory-developed tests may not have an individual FDA, CDSCO, EU IVDR, or MHRA approval date.
\clearpage
